\documentclass[10pt,a4paper]{amsart}
\usepackage[margin=19mm]{geometry}
\usepackage{amsmath,amssymb,mathtools}
\usepackage[scr=rsfs,cal=euler]{mathalfa}
\usepackage{xfrac}
\usepackage[unicode,hidelinks,bookmarks=false]{hyperref}
\newcommand{\ie}{i.e.\ }
\newcommand{\cf}{cf.\ }
\newcommand{\resp}{respectively\ }
\newcommand{\Subsection}{\subsection}
\providecommand{\nobreakdash}{\nobreak\hbox{-}}
\setlength{\emergencystretch}{2em}
\multlinegap0pt
\usepackage{bm}
\usepackage{bbm}
\usepackage{dsfont}
\usepackage{xspace}
\usepackage{cancel}
\usepackage{mathtools}
\usepackage{amsthm}
\makeatletter
\@ifundefined{lemma}{\newtheorem{lemma}{Lemma}}{}
\makeatother
\newcommand{\CC}{\mathbb{C}}
\newcommand{\PP}{\mathbb{P}}
\title[Squared Linear Models]{Squared Linear Models}
\author{Hannah Friedman, Bernd Sturmfels, Maximilian Wiesmann}
\date{Source: arXiv:2505.19351v2 (CC BY 4.0)}
\begin{document}
\maketitle

\noindent\emph{Source and licence.} Squared Linear Models. Version arXiv:2505.19351v2 (CC BY 4.0), distributed under the Creative Commons Attribution 4.0 International licence, as recorded in the arXiv licence metadata for this exact version and shown on the abstract page https://arxiv.org/abs/2505.19351v2. Full rights notice: licenses/arxiv-2505.19351-CC-BY-4.0.txt.\par\smallskip

\noindent\textit{Editorial scope.} Counted unit: the Lemma whose source label is \texttt{lem:generic\_quadric} in the article (source0.tex lines 341--345, source label \texttt{lem:generic\_quadric}), delivered together with the article's own complete original proof of it (source0.tex lines 347--390, one \texttt{proof} environment of 44 lines). No mathematics is written, repaired or re-derived: statement and proof are the article's own text line for line. Required context retained uncounted: none. Every definition, display and label of those lines is retained unchanged. Omitted from this package and not redistributed: 1--340, 346--346, 391--1909. Nothing inside a retained excerpt is omitted or altered: every retained body is delivered exactly as the article's lines are written, so no omission receipt of any kind accompanies this package. Citations: the retained excerpts carry no citation command at all, so no bibliography entry of the article is transcribed and no reference is redistributed. Reference adaptation: none was needed and none was made; every reference inside the delivered unit resolves inside it, so no unresolved reference or citation is produced. Printed slips kept literal: no printed word, symbol, hypothesis or number of the counted statement or of its proof was altered. Layout and class: the article's own class is \texttt{siamonline250211} with packages including lipsum, amsfonts, graphicx, epstopdf, algorithmic, mathtools, tikz, booktabs, tabularx, hyperref, color, enumitem, adjustbox, xspace; the wrapper is the lane's pinned \texttt{amsart} editorial wrapper at 10pt on A4, which restates the theorem-like environments the retained text needs and adds the article's own macro definitions that the retained text needs (\texttt{\textbackslash CC}, \texttt{\textbackslash PP}), with extra wrapper packages (bm, bbm, dsfont, xspace, cancel, mathtools). The delivered numbering is the unit's own. Assets: no retained span contains a figure environment, a graphics-inclusion command, a PDF-page inclusion, a table or a TikZ picture; no image file is copied into this package, the package declares no asset dependency and no image file is redistributed with it. Licence: version arXiv:2505.19351v2 of this article is distributed under the Creative Commons Attribution 4.0 International licence, as recorded in the arXiv licence metadata for this exact version and shown on the abstract page https://arxiv.org/abs/2505.19351v2. A full rights notice is delivered at licenses/arxiv-2505.19351-CC-BY-4.0.txt. Characters: every retained excerpt is pure ASCII, so the delivered engine, XeLaTeX with its default Latin Modern text font, reports no missing character.
\begin{lemma}\label{lem:generic_quadric}
    The quadratic hypersurface defined by
     $q = \ell_1^2 + \cdots + \ell_n^2$ in $\PP^{d-1}$ is in general position with respect to the hyperplane arrangement 
     $\mathcal A$ that underlies the squared linear model.
     \end{lemma}

\begin{proof}
  We prove that, for any nonempty subset $I \subseteq [n]$, the intersection $ \bigcap_{i \in I} V(\ell_i)$
    is  not tangent  to the hypersurface $V(q)$.
  Let $A $ be the $n \times d$ matrix satisfying $Ax = (\ell_i(x))_{i \in [n]}$.
  We assume that $A$ has rank $d$, so the  arrangement $\mathcal A$ is essential. 
  Otherwise, $d - {\rm rank}(A)$ of the variables can be eliminated, resulting in an essential arrangement in $\PP^{\,{\rm rank}(A) - 1}$. 
    If $A_i$ denotes the $i$th row of $A$, then the gradient of the partition function is 
  \begin{align*}
    \nabla (q(x)) \,\,= \,\,2\ell_1(x) A_1 + \cdots + 2\ell_n(x) A_n     \,\, = \,\,2x^T A^T A.
  \end{align*}
  We proceed by induction on  the cardinality $|I|$.
    We first claim that $V(\ell_i)$ is not tangent to $V(q)$ for $i = 1, \ldots, n$.
  This is equivalent to proving that the $2 \times d$ matrix
    \begin{align*}
    C \,\,= \,\,\begin{pmatrix}
      x^T A^T A\\
      A_i
    \end{pmatrix}
  \end{align*}
    has rank $2$ at every point $x \in V(q) \cap V(\ell_i)$. We now show the contrapositive.
    
  Suppose there exists  $\alpha \in \CC^*$ such that $\alpha x^T A^T A = A_i$.
  Since $A_i$ is real and nonzero and $A^TA$ is a real invertible matrix, $\alpha x$ is also real and nonzero.
  Since $q$ is a sum of squares and $\alpha \neq 0$, we have $q(\alpha x) = \alpha^2 q(x) \neq 0$.
  Thus the matrix $C$ has rank 2 for any $x \in V(q) \cap V(\ell_i)$.

We now consider an index set $I$ of cardinality $r$. After relabeling, 
$I = \{1,2,\ldots,r\}$, and we assume that  $\bigcap_{i=1}^r V(\ell_i)$  is non-empty.
We must show that this intersection is not  tangent to $V(q)$.
If it were tangent, there would exist $\lambda_1, \ldots, \lambda_r$ such that $\lambda_1 A_1 + \cdots + \lambda_r A_r =  x^TA^TA$.
  Letting ${\rm proj}_{\ell_i}(v)$ denote the orthogonal projection of $v$ onto the hyperplane $V(\ell_i)$, this implies
  \begin{align}\label{eq:tangentprojection}
  0 + \lambda_2 {\rm proj}_{\ell_1}A_2 + \cdots + \lambda_r {\rm proj}_{\ell_1}A_r 
  \,\,=\,\,  {\rm proj}_{\ell_1}(x^TA^TA).
  \end{align}
   We now view the intersections $V(\ell_2) \cap V(\ell_1), \ldots, \,V(\ell_n) \cap V(\ell_1), \,V(q) \cap V(\ell_1)$ as subvarieties of $V(\ell_1) \cong \PP^{d - 2}$.
  In that ambient space, the normal vector of $V(q) \cap V(\ell_1)$ is ${\rm proj}_{\ell_1}(x^TA^TA)$ at $x$.
  Similarly,  the normal vector of $V(\ell_i) \cap V(\ell_1)$ is ${\rm proj}_{\ell_1}(A_i)$.
  The condition  \eqref{eq:tangentprojection} means that the intersection 
  $\bigcap_{i=2}^r(V(\ell_i) \cap V(\ell_1))$  is tangent to 
  $V(q) \cap V(\ell_1) = V(\ell_2^2 + \cdots + \ell_n^2) \cap V(\ell_1)$
  inside $ \PP^{d - 2}$.
  This is a contradiction to our induction hypothesis. We conclude that the $\ell_i$ are in general position relative to the quadric $q$ in the complex projective space $\PP^{d-1}$.
\end{proof}

\end{document}
